\documentclass[12pt]{article}
\usepackage[margin=1in]{geometry}
\usepackage{amsmath,amsfonts,mathtools,amsthm,amssymb}
\usepackage[l2tabu,orthodox]{nag}
\usepackage{microtype} % fixes spacing or whatever
\usepackage{enumitem}
\usepackage{tikz-cd}
\usepackage{xcolor}
\usepackage{physics}
\usepackage{mathrsfs}
\usepackage{cancel}
% \usepackage{libertine,libertinust1math}
\usepackage{eucal}
\usepackage[toc,page]{appendix}
\usepackage{pgfplots}
\pgfplotsset{compat=1.18}

\newtheorem{proposition}{Proposition}
\newtheorem{theorem}{Theorem}
\newtheorem{lemma}{Lemma}
\newtheorem{corollary}{Corollary}

\theoremstyle{definition}
\newtheorem{remark}{Remark}
\newtheorem{definition}{Definition}
\newtheorem{example}{Example}

% \usepackage[sc,osf]{mathpazo}   % With old-style figures and real smallcaps.
% \linespread{1.025}              % Palatino leads a little more leading
% % Euler for math and numbers
% \usepackage[euler-digits,small]{eulervm}
% \AtBeginDocument{\renewcommand{\hbar}{\hslash}}

\usepackage{etoolbox}
\usepackage{dynkin-diagrams}
\def\row#1/#2!{\dynkin{#1}{#2}\\}
\newcommand{\tble}[1]{
   \renewcommand*\do[1]{\row##1!}
   \[
      \begin{array}{ll}\docsvlist{#1}\end{array}
   \]
}

\allowdisplaybreaks

\renewcommand{\epsilon}{\varepsilon}
% \renewcommand{\phi}{\varphi}

\DeclareMathAlphabet{\mathpzc}{OT1}{pzc}{m}{it}

\newcommand{\goth}[1]{\mathfrak{#1}}
\newcommand{\red}[1]{#1_{\text{red}}}
\newcommand{\ad}[1]{#1_{\text{Ad}}}
\newcommand{\reg}[1]{#1_{\text{reg}}}
\newcommand{\sh}[1]{#1^{\text{sh}}}
\newcommand{\cpdv}[2]{\cfrac{\partial #1}{\partial #2}}

\newcommand{\fA}{\mathcal{A}}
\newcommand{\fB}{\mathcal{B}}
\newcommand{\fC}{\mathcal{C}}
\newcommand{\fD}{\mathcal{D}}
\newcommand{\fE}{\mathcal{E}}
\newcommand{\fF}{\mathcal{F}}
\newcommand{\fG}{\mathcal{G}}
\newcommand{\fH}{\mathcal{H}}
\newcommand{\fI}{\mathcal{I}}
\newcommand{\fJ}{\mathcal{J}}
\newcommand{\fK}{\mathcal{K}}
\newcommand{\fL}{\mathcal{L}}
\newcommand{\fM}{\mathcal{M}}
\newcommand{\fN}{\mathcal{N}}
\newcommand{\fO}{\mathcal{O}}
\newcommand{\fP}{\mathcal{P}}
\newcommand{\fQ}{\mathcal{Q}}
\newcommand{\fR}{\mathcal{R}}
\newcommand{\fS}{\mathcal{S}}
\newcommand{\fT}{\mathcal{T}}
\newcommand{\fX}{\mathcal{X}}
\newcommand{\fY}{\mathcal{Y}}
\newcommand{\fZ}{\mathcal{Z}}

\newcommand{\be}{\mathbf{e}}
\newcommand{\bx}{\mathbf{x}}
\newcommand{\bone}{\mathbf{1}}
\newcommand{\cx}{\bar{x}}
\newcommand{\cy}{\bar{y}}
\newcommand{\vx}{\vec{x}}
\newcommand{\vsigma}{\vec{\sigma}}
\newcommand{\tzeta}{\tilde{\zeta}}

\newcommand{\gm}{\goth{m}}
\newcommand{\gp}{\goth{p}}
\newcommand{\gF}{\goth{F}}
\newcommand{\gU}{\goth{U}}
\newcommand{\gV}{\goth{V}}

\newcommand{\A}{\mathbf{A}}
\newcommand{\C}{\mathbf{C}}
\newcommand{\F}{\mathbf{F}}
\newcommand{\G}{\mathbf{G}}
\newcommand{\bH}{\mathbf{H}}
\newcommand{\N}{\mathbf{N}}
\newcommand{\PP}{\mathbf{P}}
\newcommand{\Q}{\mathbf{Q}}
\newcommand{\R}{\mathbf{R}}
\newcommand{\Z}{\mathbf{Z}}

\newcommand{\dlog}{\dd\log}

\newcommand\srestr[2]{{\left.\kern-\nulldelimiterspace #1\vphantom{\small|} \right|_{#2}}}
\newcommand\restr[2]{{\left.\kern-\nulldelimiterspace #1 \vphantom{\big|} \right|_{#2}}}
\newcommand\gsec[2]{\Gamma({#1}, {#2})}
\newcommand\gsecx[1]{\Gamma(X, {#1})}

\DeclareMathOperator{\chr}{char}
\DeclareMathOperator{\rProj}{\mathbf{Proj}}
\DeclareMathOperator{\rSpec}{\mathbf{Spec}}
\DeclareMathOperator{\rHom}{\mathbf{Hom}}
\DeclareMathOperator{\rExt}{\mathbf{Ext}}
\DeclareMathOperator{\id}{id}
\DeclareMathOperator{\Frac}{Frac}
\DeclareMathOperator{\rk}{rank}
\DeclareMathOperator{\deter}{det}
\DeclareMathOperator{\pic}{Pic}
\DeclareMathOperator{\cacl}{CaCl}
\DeclareMathOperator{\trd}{tr.d.}
\DeclareMathOperator{\cl}{Cl}
\DeclareMathOperator{\depth}{depth}
\DeclareMathOperator{\codim}{codim}
\DeclareMathOperator{\Div}{Div}
\DeclareMathOperator{\coker}{coker}
\DeclareMathOperator{\len}{length}
\DeclareMathOperator{\height}{height}
\DeclareMathOperator{\supp}{Supp}
\DeclareMathOperator{\proj}{Proj}
\DeclareMathOperator{\im}{im}
\DeclareMathOperator{\Hom}{Hom}
\DeclareMathOperator{\Der}{Der}
\DeclareMathOperator{\spec}{Spec}
\DeclareMathOperator{\Aut}{Aut}
\DeclareMathOperator{\ch}{char}
\DeclareMathOperator{\tor}{Tor}
\DeclareMathOperator{\Ann}{Ann}
\DeclareMathOperator{\Syl}{Syl}
\DeclareMathOperator{\Sym}{Sym}
\DeclareMathOperator{\GL}{GL}
\DeclareMathOperator{\SL}{SL}
\DeclareMathOperator{\Stab}{Stab}
\DeclareMathOperator{\Perm}{Perm}
\DeclareMathOperator{\Orb}{Orb}
\DeclareMathOperator{\Gal}{Gal}
\DeclareMathOperator{\Supp}{Supp}
\DeclareMathOperator{\Ext}{Ext}
\DeclareMathOperator{\Tor}{Tor}
\DeclareMathOperator{\PGL}{PGL}
\DeclareMathOperator{\hd}{hd}
\DeclareMathOperator{\pd}{pd}
\DeclareMathOperator{\SO}{SO}
\DeclareMathOperator{\SU}{SU}
\DeclareMathOperator{\Sp}{Sp}
\DeclareMathOperator{\Oth}{O}
\DeclareMathOperator{\Uni}{U}
\DeclareMathOperator{\evf}{ev}
\DeclareMathOperator{\cH}{H}
\DeclareMathOperator{\dR}{R}
\DeclareMathOperator{\End}{End}
\DeclareMathOperator{\Hasse}{Hasse}
\DeclareMathOperator{\Num}{Num}

\author{James Lee}

% Solarized
% \pagecolor[RGB]{0,20,26}
% \color[RGB]{191,191,191}

% Sephia
% \pagecolor[RGB]{249,239,220}

% Grey
% \pagecolor[RGB]{200, 200, 200}

% Black
% \pagecolor[RGB]{0,0,0}
% \color[RGB]{255,255,255}

\title{Chapter 4, Section 5}

\begin{document}
\maketitle
\begin{enumerate} [label=\textbf{\arabic*.}, leftmargin=0em]

  \item Show that a hyperelliptic curve can never be a complete intersection in any projective space.

        \begin{proof}
          Combine (Ex. 3.3) and (5.2).
        \end{proof}

  \item If $X$ is a curve of genus $\geq 2$ over a field of characteristic 0, show that the group $\Aut{X}$ of automorphisms of $X$ is finite.

        \begin{proof}
          Suppose $X$ is hyperelliptic, and let $f : X \to \PP^1$ be a 2-fold covering $f : X \to \PP^1$ given by the unique $g_2^1$.
          The ramification locus consists of $2g + 2$ distinct points.
          After a suitable automorphism of $\PP^1$, we can label them $0, 1, \infty, \beta_1, \dots, \beta_{2g - 1}$.
          Let $\sigma$ be an automorphism of $X$.
          Then $f \circ \sigma$ gives a $g_2^1$, so by uniqueness of $g_2^1$, there exists an automorphism $\tau$ of $\PP^1$ such that the following diagram commutes
          \[ \begin{tikzcd}
              X \arrow[d] \arrow[r, "\sigma"] & X \arrow[d] \\
              \PP^1 \arrow[r, "\tau"]         & \PP^1
            \end{tikzcd}. \]
            Thus, $\sigma$ permutes the ramification points of $f$ which induces a group homomorphism $\Aut{X} \to \Sigma_{2g - 1}$.
            If $\sigma$ is in the kernel, i.e., $\sigma$ fixes the ramification points of $f$,

            If $X$ is not hyperelliptic, we show that $\Aut(X)$ permutes the hyperosculation points of the canonical embedding.
        \end{proof}

  \item \textit{Moduli of Curves of Genus 4.} The hyperelliptic curves of genus 4 form an irreducible family of dimension 7.
        The nonhyperelliptic ones form an irreducible family of dimension 9.
        The subset of those having only one $g_3^1$ is an irreducible family of dimension 8.

        \begin{proof}
          Hyperelliptic elliptic curves of genus 4 are determined as two-fold coverings of $\PP^1$, ramified at $0, 1, \infty$, and $2g - 1 = 7$ additional points, up to the action of a certain finite group.
          If $X$ is a nonhyperelliptic curve of genus 4, then its canonical embedding in $\PP^3$ is the complete intersection of a unique irreducible quadric surface and an irreducible cubic surface.
          Conversely, if $X$ is a nonsingular curve in $\PP^3$ which is a complete intersection of a quadric and a cubic surface, then $X$ is a canonical curve of genus 4.
          There are only two isomorphism classes of irreducible quadric surfaces in $\PP^3$, so the family of all nonhyperelliptic curves of genus 4 is parametrized by two copies of an open set $U \subseteq \PP^{N}$ with $N = 19$, because a form of degree 3 has 20 coefficients.
          Fixing a quadric surface $Q \subseteq \PP^3$, there is a dominant morphism $U \to \goth{M}_4$ defined by $F \mapsto Q \cap F$, so it remains to compute the dimension of a general fibre.
          Let $F, F'$ be two irreducible cubic surfaces.
          Since the embedding is canonical, $Q \cap F$ and $Q \cap F'$ are  isomorphic if and only if they differ by an automorphism of $\PP^3$.
        \end{proof}

  \item[\textbf{5.}] \textit{Curves of Genus 5.} Assume $X$ is not hyperelliptic.
        \begin{itemize}
          \item[(a)] The curves of genus 5 whose canonical model in $\PP^4$ is a complete intersection $F_1.F_2.F_3$ form a family of dimension 12.
          \item[(b)] $X$ has a $g_3^1$ if and only if it can be represented as a plane quintic with one node.
                These form an irreducible family of dimension 11.
          \item[(c)] In that case, the conics through the node cut out the canonical system (not counting the fixed points at the node).
                Mapping $\PP^2 \to \PP^4$ by this linear system of conics, show that the canonical curve $X$ is contained in a cubic surface $V \subseteq \PP^4$, with $V$ isomorphic to $\PP^2$ with one point blown up (II, Ex. 7.7).
                Furthermore, $V$ is the union of all the trisecants of $X$ corresponding to the $g_3^1$ (5.5.3), so $V$ is contained in the intersection of all the quadric hypersurfaces containing $X$.
                Thus $V$ and $g_3^1$ are unique.
        \end{itemize}

        \begin{proof} $ $ \vspace{0pt}
          \begin{itemize} [leftmargin=0cm]
            \item[(a)]

            \item[(b)]

            \item[(c)]
          \end{itemize}
        \end{proof}

  \item[\textbf{7.}] \begin{itemize}
          \item[(a)] Any automorphism of a curve of genus 3 is induced by an automorphism of $\PP^2$ via the canonical embedding.
          \item[(b)] Assume $\chr{k} \neq 3$.
                If $X$ is the curve given by
                \begin{equation*}
                  x^3y + y^3z + z^3 x = 0,
                \end{equation*}
                the group $\Aut{X}$ is the simple group of order 168, whose order is the maximum $84(g - 1)$ allowed by (Ex. 2.5).
          \item[(c)] Most curves of genus 3 have no automorphisms except the identity.
        \end{itemize}

        \begin{proof} $ $ \vspace{0pt}
          \begin{itemize} [leftmargin=0cm]
            \item[(a)]

            \item[(b)]

            \item[(c)]
          \end{itemize}
        \end{proof}

\end{enumerate}

\end{document}
